\documentclass[11pt,a4paper]{article}

\usepackage[utf8]{inputenc}
\usepackage[
  left=1.05in,
  right=1.05in,
  top=1in,
  bottom=1in,
  headheight=14pt
]{geometry}
\usepackage{amsmath,amssymb,amsthm,mathtools}
\usepackage{bm}
\usepackage{enumitem}
\usepackage{booktabs,array}
\usepackage{microtype}
\usepackage{xcolor}
\usepackage[breakable,skins]{tcolorbox}
\usepackage[hidelinks]{hyperref}
\usepackage{fancyhdr}
% Start a new page unless #1 of vertical space remains.
\makeatletter
\newcommand{\Needspace}[1]{\par\penalty-100\begingroup
  \dimen@=#1\relax
  \dimen@ii=\pagegoal\advance\dimen@ii by -\pagetotal
  \ifdim\dimen@>\dimen@ii\ifdim\dimen@ii>\z@\vfil\fi\break\fi
  \endgroup}
\makeatother

\linespread{1.25}
\setlength{\parindent}{0pt}
\setlength{\parskip}{1em plus 0.25em minus 0.1em}
\setlength{\abovedisplayskip}{1.2em plus 0.3em minus 0.2em}
\setlength{\belowdisplayskip}{1.2em plus 0.3em minus 0.2em}
\setlength{\abovedisplayshortskip}{0.9em plus 0.2em}
\setlength{\belowdisplayshortskip}{1em plus 0.2em}
\setlength{\jot}{9pt}
\setlength{\emergencystretch}{2em}
\raggedbottom

\setlist[itemize]{leftmargin=*,itemsep=6pt,topsep=6pt,parsep=3pt}
\setlist[enumerate]{leftmargin=*,itemsep=6pt,topsep=6pt,parsep=3pt}

% Plain black frame for the question statement.
\newtcolorbox{qbox}{
  enhanced,
  colback=white,colframe=black,
  boxrule=0.6pt,sharp corners,
  left=10pt,right=10pt,top=4pt,bottom=6pt,
  before skip=14pt,after skip=10pt
}

\newcounter{qn}
% \Q{title}{type}{marks}
\newcommand{\Q}[3]{\Needspace{14\baselineskip}\refstepcounter{qn}\subsection*{Q\theqn.\ #1}%
  \addcontentsline{toc}{subsection}{Q\theqn.\ #1}%
  \noindent\textit{#2}\hfill\textit{[#3 marks]}\par}

\newenvironment{solution}
  {\par\noindent\textbf{Solution.}\par\nobreak}
  {\par\hfill\(\square\)\par}

\newcommand{\step}[1]{\par\addvspace{0.5em}\noindent\textbf{#1}\par\nobreak}
\newcommand{\pd}[2]{\frac{\partial #1}{\partial #2}}
\newcommand{\pdd}[2]{\frac{\partial^2 #1}{\partial #2^2}}
\newcommand{\hb}{\bar h}

\pagestyle{fancy}
\fancyhf{}
\fancyhead[L]{MTL712}
\fancyhead[R]{Practice: Notes 4 and 5}
\fancyfoot[C]{\thepage}
\renewcommand{\headrulewidth}{0.4pt}
\renewcommand{\footrulewidth}{0pt}

\title{\textbf{MTL712: Computational Methods for Differential\\Equations}\\[6pt]
\large Practice Questions with Complete Solutions\\
on Notes 4 and 5}
\author{\normalsize Multistep methods, BVPs, parabolic PDEs, ADI and
three-level schemes}
\date{}

\begin{document}

\maketitle
\thispagestyle{plain}

These questions cover every topic of Notes 4 and 5 and are written in
the style of the major examinations: proofs, derivations, numerical
calculations, stability analyses and short conceptual answers.  Each
question states its type and suggested marks.  Try each question before
reading its solution.  All numerical answers have been checked by
computer.

\begin{center}
\renewcommand{\arraystretch}{1.25}
\begin{tabular}{@{}lll@{}}
\toprule
Section & Topic & Questions\\
\midrule
1 & Linear multistep methods for \(y'=f\) & Q1--Q7\\
2 & Methods for \(y''=f\) & Q8--Q10\\
3 & Boundary value problems & Q11--Q15\\
4 & Classification and difference operators & Q16--Q17\\
5 & Parabolic equations: schemes, LTE, stability & Q18--Q27\\
6 & Lecture 35: ADI, three-level and Du Fort--Frankel & Q28--Q34\\
\bottomrule
\end{tabular}
\end{center}

{\linespread{1.0}\selectfont\setlength{\parskip}{0pt}\tableofcontents}
\clearpage

%#####################################################################
\Needspace{32\baselineskip}
\section{Linear Multistep Methods for \(y'=f(x,y)\)}
%#####################################################################

Notation as in Note 4: \(\rho(\zeta)=\sum\alpha_j\zeta^j\),
\(\sigma(\zeta)=\sum\beta_j\zeta^j\), \(w=\zeta-1\), and the error
constant of an order-\(p\) method is \(C_{p+1}\).

%---------------------------------------------------------------------
\Q{Two methods from \(\rho(\zeta)=\zeta^2-1\)}{Numerical / Henrici's theorem}{10}
\begin{qbox}
Let \(\rho(\zeta)=\zeta^2-1\).
\begin{enumerate}[label=(\alph*)]
  \item Use Henrici's construction to find the \emph{explicit} method of
    maximal order.  State its order and error constant.
  \item Find the \emph{implicit} method of maximal order, its order and
    error constant.
  \item Show that the method in (a) has no interval of absolute
    stability.
\end{enumerate}
\end{qbox}

\begin{solution}
\step{Expansion}
With \(w=\zeta-1\), \(\rho=w(2+w)\), and using
\(\frac{w}{\log(1+w)}=1+\frac w2-\frac{w^2}{12}+\frac{w^3}{24}-\frac{19w^4}{720}+\cdots\),
\[
  \frac{\rho(\zeta)}{\log\zeta}
  =(2+w)\Bigl[1+\frac w2-\frac{w^2}{12}+\frac{w^3}{24}-\frac{19w^4}{720}+\cdots\Bigr]
  =2+2w+\frac13w^2+0\cdot w^3-\frac1{90}w^4+\cdots .
\]

\step{(a) Explicit (\(k'=k-1=1\))}
Keep terms up to \(w\): \(\sigma=2+2w=2\zeta\).  The method is the
\emph{mid-point rule}
\[
  y_{n+2}-y_n=2hf_{n+1}.
\]
The first omitted term is \(\frac13w^2\), so the order is \(2\) and the
error constant is \(C_3=\frac13\).

\step{(b) Implicit (\(k'=k=2\))}
Keep terms up to \(w^2\):
\(\sigma=2+2w+\frac13w^2=\frac13(\zeta^2+4\zeta+1)\).  This is
\emph{Milne--Simpson},
\[
  y_{n+2}-y_n=\frac h3\bigl(f_{n+2}+4f_{n+1}+f_n\bigr).
\]
The \(w^3\) coefficient vanishes, so the order is \(4\), with
\(C_5=-\frac1{90}\).

\step{(c) Absolute stability of the mid-point rule}
The stability polynomial is \(\pi(z,\hb)=z^2-2\hb z-1\), with roots
\(z_{1,2}=\hb\pm\sqrt{\hb^2+1}\).  Their product is \(-1\), so
\(|z_1||z_2|=1\), and both cannot be less than \(1\) in modulus.  Indeed
for \(\hb>0\), \(z_1=\hb+\sqrt{\hb^2+1}>1\); for \(\hb<0\),
\(z_2=\hb-\sqrt{\hb^2+1}<-1\); and for \(\hb=0\) the roots are \(\pm1\).
Hence the method is absolutely stable for no real \(\hb\): the interval
of absolute stability is empty.  (It is zero-stable, since the roots
\(\pm1\) of \(\rho\) are simple, but only weakly stable.)
\end{solution}

%---------------------------------------------------------------------
\Q{Simpson's three-eighths method}{Numerical}{8}
\begin{qbox}
For the three-step method
\[
  y_{n+3}-y_n=\frac{3h}8\bigl(f_{n+3}+3f_{n+2}+3f_{n+1}+f_n\bigr),
\]
\begin{enumerate}[label=(\alph*)]
  \item show that it is consistent and zero-stable;
  \item find its order and error constant;
  \item explain why its order is the largest possible for a zero-stable
    three-step method.
\end{enumerate}
\end{qbox}

\begin{solution}
Here \(\alpha=(-1,0,0,1)\), \(\beta=\frac38(1,3,3,1)\),
\(\rho(\zeta)=\zeta^3-1\), \(\sigma(\zeta)=\frac38(\zeta^3+3\zeta^2+3\zeta+1)\).

\step{(a)}
\(\rho(1)=0\) and \(\rho'(1)=3=\sigma(1)=\frac38\cdot8\), so the method
is consistent.  The roots of \(\rho\) are the cube roots of unity
\(1,e^{\pm2\pi i/3}\): all have modulus \(1\) and are simple, so the
root condition holds and the method is zero-stable (weakly stable,
since roots other than \(1\) lie on the unit circle).

\step{(b)}
Use \(C_q=\frac1{q!}\sum j^q\alpha_j-\frac1{(q-1)!}\sum j^{q-1}\beta_j\).
First \(C_0=\rho(1)=0\) and \(C_1=\rho'(1)-\sigma(1)=3-3=0\).  For
\(q\ge2\), computing each sum directly,
\[
  \sum_jj^q\alpha_j=3^q,
  \qquad
  \sum_jj^{q-1}\beta_j=\tfrac38\bigl(3^{q-1}+3\cdot2^{q-1}+3\bigr)\quad(q\ge2),
\]
and
\begin{align*}
  C_2&=\tfrac92-\tfrac38(3+6+3)=\tfrac92-\tfrac92=0,\\
  C_3&=\tfrac{27}6-\tfrac38\cdot\tfrac{9+12+3}{2}=\tfrac92-\tfrac92=0,\\
  C_4&=\tfrac{81}{24}-\tfrac38\cdot\tfrac{27+24+3}{6}=\tfrac{27}8-\tfrac{27}8=0,\\
  C_5&=\tfrac{243}{120}-\tfrac38\cdot\tfrac{81+48+3}{24}
      =\tfrac{81}{40}-\tfrac{33}{16}=-\tfrac3{80}\ne0 .
\end{align*}
So the order is \(4\) and the error constant is \(C_5=-\frac3{80}\).

\step{(c)}
By the first Dahlquist barrier, a zero-stable \(k\)-step method has
order at most \(k+1\) when \(k\) is odd.  For \(k=3\) the maximum is
\(4\), which this method attains.
\end{solution}

%---------------------------------------------------------------------
\Q{BDF2 from \(\sigma(\zeta)=\zeta^2\)}{Derivation}{8}
\begin{qbox}
Find the polynomial \(\rho\) of degree \(2\) such that the method with
\(\sigma(\zeta)=\zeta^2\) has maximal order.  Write the method, find its
order and error constant, and check zero-stability.  Also give the
method normalised so that \(\alpha_2=1\).
\end{qbox}

\begin{solution}
\step{Construction}
For given \(\sigma\), expand \(\log\zeta\cdot\sigma(\zeta)\) in powers of
\(w\) and truncate at degree \(k=2\):
\[
  \log(1+w)\,(1+w)^2
  =\Bigl(w-\frac{w^2}2+\frac{w^3}3-\cdots\Bigr)(1+2w+w^2)
  =w+\frac32w^2+\frac13w^3+\cdots .
\]
So
\[
  \rho=w+\frac32w^2=(\zeta-1)+\frac32(\zeta-1)^2
  =\frac32\zeta^2-2\zeta+\frac12,
\]
and the method is
\[
  \frac32y_{n+2}-2y_{n+1}+\frac12y_n=hf_{n+2}
  \qquad\text{(BDF2)}.
\]

\step{Order and error constant}
\[
  \varphi(\zeta)=\frac{\rho-\log\zeta\cdot\sigma}{\log\zeta}
  =\frac{-\frac13w^3+\cdots}{w+\cdots}=-\frac13w^2+\cdots,
\]
so by the order lemma the order is \(2\) and \(C_3=-\frac13\).
Normalised (divide by \(\frac32\)):
\[
  y_{n+2}-\frac43y_{n+1}+\frac13y_n=\frac23hf_{n+2},
  \qquad C_3=-\frac29 .
\]

\step{Zero-stability}
\(\rho(\zeta)=\frac12(3\zeta-1)(\zeta-1)\) has roots \(1\) and
\(\frac13\): the root condition holds, and \(1\) is the only root on the
unit circle, so the method is strongly stable.  Consistency:
\(\rho(1)=0\), \(\rho'(1)=3-2=1=\sigma(1)\).
\end{solution}

%---------------------------------------------------------------------
\Q{A convergent method is consistent}{Proof}{8}
\begin{qbox}
Prove that a convergent linear multistep method satisfies
\(\rho(1)=0\) and \(\rho'(1)=\sigma(1)\).  You may assume that a
convergent method is zero-stable.
\end{qbox}

\begin{solution}
\step{\(\rho(1)=0\)}
Apply the method to \(y'=0\), \(y(0)=1\) (solution \(y\equiv1\)), with
exact starting values \(y_0=\dots=y_{k-1}=1\).  The difference equation
\(\sum_j\alpha_jy_{n+j}=0\) does not involve \(h\), so \(y_n\) does not
depend on \(h\), and convergence (\(y_n\to1\) as \(h\to0\) with
\(nh=x\) fixed, i.e.\ \(n\to\infty\)) means \(\lim_{n\to\infty}y_n=1\).
Letting \(n\to\infty\) in the difference equation gives
\(\sum_j\alpha_j=\rho(1)=0\).

\step{\(\rho'(1)=\sigma(1)\)}
Apply the method to \(y'=1\), \(y(0)=0\) (solution \(y=x\)):
\[
  \sum_j\alpha_jy_{n+j}=h\,\sigma(1).
\]
Since \(\rho(1)=0\) and the method is zero-stable, \(1\) is a simple root
of \(\rho\), so \(\rho'(1)\ne0\).  Put \(K=\sigma(1)/\rho'(1)\) and
\(y_n=Knh\).  Then
\[
  \sum_j\alpha_jK(n+j)h=Kh\bigl(n\rho(1)+\rho'(1)\bigr)=Kh\rho'(1)=h\sigma(1),
\]
so \(y_n\) solves the difference equation, with starting values
\(Ksh\to0\).  Convergence forces \(y_n=Knh=Kx\to x\) for every \(x>0\),
hence \(K=1\), i.e.\ \(\rho'(1)=\sigma(1)\).
\end{solution}

%---------------------------------------------------------------------
\Q{Absolute instability for small positive \(\hb\)}{Proof}{10}
\begin{qbox}
Prove that every consistent zero-stable LMM is absolutely unstable for
all sufficiently small \(\hb>0\).  Explain why this is not a defect.
\end{qbox}

\begin{solution}
Let the order be \(p\ge1\).  Expanding the exponentials,
\[
  \pi(e^{\hb},\hb)=\rho(e^{\hb})-\hb\,\sigma(e^{\hb})
  =\sum_{q\ge0}C_q\hb^{\,q}=O(\hb^{\,p+1}).
\]
Factorise \(\pi(z,\hb)=(\alpha_k-\hb\beta_k)\prod_{s=1}^k(z-r_s(\hb))\).
As \(\hb\to0\), \(r_s(\hb)\to\zeta_s\), the roots of \(\rho\).
Consistency gives \(\rho(1)=0\) and zero-stability makes this root
simple, say \(\zeta_1=1\), \(\zeta_s\ne1\) for \(s\ge2\).  So in
\[
  \pi(e^{\hb},\hb)=(\alpha_k-\hb\beta_k)\prod_s\bigl(e^{\hb}-r_s(\hb)\bigr)
\]
every factor except \(e^{\hb}-r_1(\hb)\) tends to a non-zero limit.
Hence \(|e^{\hb}-r_1(\hb)|\le K\hb^{\,p+1}\), and
\[
  |r_1(\hb)|\ge e^{\hb}-K\hb^{\,p+1}\ge1+\hb-K\hb^{\,p+1}>1
\]
once \(K\hb^{\,p}<1\).  So some root has modulus greater than \(1\): the
method is absolutely unstable for small \(\hb>0\).

This is not a defect: for \(\lambda>0\) the exact solution
\(e^{\lambda x}\) grows, and \(r_1\approx e^{\hb}\) is the root that
correctly reproduces this growth.
\end{solution}

%---------------------------------------------------------------------
\Q{Short answers on multistep methods}{Short answer}{8}
\begin{qbox}
Answer briefly, with reasons.
\begin{enumerate}[label=(\alph*)]
  \item Can a zero-stable three-step method have order \(5\)?
  \item Is the method
    \(y_{n+2}-3y_{n+1}+2y_n=h(\beta_2f_{n+2}+\beta_1f_{n+1}+\beta_0f_n)\)
    convergent for some choice of the \(\beta_j\)?
  \item A \(k\)-step method has \(2k+2\) coefficients.  Why can it reach
    order \(2k\) at most, and what does zero-stability cost?
  \item State Dahlquist's equivalence theorem.
\end{enumerate}
\end{qbox}

\begin{solution}
\step{(a)}
No.  By the first Dahlquist barrier the order of a zero-stable
\(k\)-step method is at most \(k+1\) for odd \(k\); for \(k=3\) that is
\(4\).

\step{(b)}
No.  \(\rho(\zeta)=(\zeta-1)(\zeta-2)\) has the root \(2\) outside the
unit disc, so the method is not zero-stable for any \(\sigma\), and
zero-stability is necessary for convergence.

\step{(c)}
After normalising (\(\alpha_k=1\)) there are \(2k+1\) free coefficients,
so we can impose \(C_0=\dots=C_{2k}=0\), giving order \(2k\).
Zero-stability reduces the attainable order to \(k+1\) (odd \(k\)) or
\(k+2\) (even \(k\)): roughly half.

\step{(d)}
An LMM applied to \(y'=f(x,y)\), with \(f\) Lipschitz in \(y\) and
consistent starting values, is convergent if and only if it is
consistent and zero-stable.
\end{solution}

%---------------------------------------------------------------------
\Q{Henrici's theorem}{Proof}{8}
\begin{qbox}
Let \(\rho\) have degree \(k\) with \(\rho(1)=0\), \(\rho'(1)\ne0\), and
let \(0\le k'\le k\).  Prove that there is a unique \(\sigma\) of degree
\(\le k'\) such that the method \((\rho,\sigma)\) has order
\(\ge k'+1\).  (You may use: the order is \(p\) iff
\(\varphi=\rho/\log\zeta-\sigma\) has a zero of order \(p\) at
\(\zeta=1\).)
\end{qbox}

\begin{solution}
\(\rho\) and \(\log\zeta\) both have a simple zero at \(\zeta=1\), so
\(\rho/\log\zeta\) is holomorphic near \(1\):
\[
  \frac{\rho(\zeta)}{\log\zeta}=\sum_{m\ge0}c_m(\zeta-1)^m .
\]
From \(\rho=\log\zeta\sum c_m(\zeta-1)^m\), differentiating and setting
\(\zeta=1\) gives \(c_0=\rho'(1)\ne0\).  Define
\(\sigma(\zeta)=\sum_{m=0}^{k'}c_m(\zeta-1)^m\).  Then
\[
  \varphi(\zeta)=\sum_{m>k'}c_m(\zeta-1)^m
\]
has a zero of order at least \(k'+1\) at \(\zeta=1\), so the order is at
least \(k'+1\).  Also \(\sigma(1)=c_0=\rho'(1)\), so the method is
consistent.  Conversely, any \(\sigma\) of degree \(\le k'\) giving
order \(\ge k'+1\) must agree with the Taylor series of
\(\rho/\log\zeta\) up to \((\zeta-1)^{k'}\); since Taylor coefficients
are unique, \(\sigma\) is unique.
\end{solution}

%#####################################################################
\Needspace{32\baselineskip}
\section{Methods for \(y''=f(x,y)\)}
%#####################################################################

%---------------------------------------------------------------------
\Q{A one-parameter family containing Störmer and Numerov}{Derivation}{8}
\begin{qbox}
Consider
\[
  y_{n+1}-2y_n+y_{n-1}=h^2\bigl[\beta f_{n+1}+(1-2\beta)f_n+\beta f_{n-1}\bigr].
\]
\begin{enumerate}[label=(\alph*)]
  \item Show that it has order at least \(2\) for every \(\beta\), and
    find the error constant.
  \item For which \(\beta\) is the order higher?  Find that order and
    the error constant.
\end{enumerate}
\end{qbox}

\begin{solution}
Expand about \(x_n\).  With
\(y(x_{n\pm1})=\sum\frac{(\pm h)^q}{q!}y^{(q)}\) and the right-hand side
weights \(\beta,1-2\beta,\beta\) at \(x_{n+1},x_n,x_{n-1}\), the method
is symmetric, so all odd \(C_q\) vanish.  Using the two-step
coefficients \(C_0,\dots,C_6\) derived for Störmer and Numerov in
Note 4, with \(a_1=-2\), \(a_2=1\), \(b_0=b_2=\beta\),
\(b_1=1-2\beta\):
\begin{align*}
  C_0&=1-2+1=0,\qquad C_1=1-1=0,\\
  C_2&=\tfrac12(1+1)-(\beta+1-2\beta+\beta)=1-1=0,\\
  C_4&=\tfrac1{24}\cdot2-\tfrac12\cdot2\beta=\tfrac1{12}-\beta,\\
  C_6&=\tfrac1{720}\cdot2-\tfrac1{24}\cdot2\beta=\tfrac1{360}-\tfrac\beta{12}.
\end{align*}

\step{(a)}
\(C_0=\dots=C_3=0\) for every \(\beta\), so the order is at least \(2\).
If \(\beta\ne\frac1{12}\), the order is exactly \(2\) with
\(C_4=\frac1{12}-\beta\).  (\(\beta=0\) gives Störmer with
\(C_4=\frac1{12}\).)

\step{(b)}
For \(\beta=\frac1{12}\), \(C_4=C_5=0\) and
\(C_6=\frac1{360}-\frac1{144}=-\frac1{240}\ne0\), so the order is \(4\).
This is Numerov's method,
\[
  y_{n+1}-2y_n+y_{n-1}=\frac{h^2}{12}\bigl(f_{n+1}+10f_n+f_{n-1}\bigr),
  \qquad C_6=-\frac1{240}.
\]
Since \(\sigma(1)=1\), the normalised error constants are the same.
\end{solution}

%---------------------------------------------------------------------
\Q{Numerov's method for \(y''=y\)}{Numerical}{8}
\begin{qbox}
Use Numerov's method with \(h=0.1\) to approximate \(y(0.2)\) and
\(y(0.3)\) for \(y''=y\), \(y(0)=1\), \(y'(0)=0\).  Obtain \(y_1\) from
a fourth-order Taylor series.  Compare with the exact solution
\(\cosh x\).
\end{qbox}

\begin{solution}
\step{Starting value}
\(y''=y\) gives \(y''(0)=1\), \(y'''(0)=y'(0)=0\), \(y^{(4)}(0)=1\), so
\[
  y_1=1+\frac{h^2}2+\frac{h^4}{24}=1.0050041667 .
\]

\step{Recurrence}
With \(f=y\), Numerov is linear in \(y_{n+1}\):
\[
  \Bigl(1-\frac{h^2}{12}\Bigr)y_{n+1}
  =\Bigl(2+\frac{10h^2}{12}\Bigr)y_n-\Bigl(1-\frac{h^2}{12}\Bigr)y_{n-1}.
\]
For \(h=0.1\): \(1-\frac{h^2}{12}=0.9991666667\) and
\(2+\frac{10h^2}{12}=2.0166666667\).

\step{Steps}
\begin{align*}
  y_2&=\frac{2.0166666667\times1.0050041667-0.9991666667\times1}{0.9991666667}
     =1.0200667570,\\
  y_3&=\frac{2.0166666667\times1.0200667570-0.9991666667\times1.0050041667}{0.9991666667}
     =1.0453385226 .
\end{align*}

\step{Comparison}
\begin{center}
\begin{tabular}{@{}cccc@{}}
\toprule
\(x\) & Numerov & \(\cosh x\) & error\\
\midrule
0.2 & 1.0200667570 & 1.0200667556 & \(1.4\times10^{-9}\)\\
0.3 & 1.0453385226 & 1.0453385141 & \(8.5\times10^{-9}\)\\
\bottomrule
\end{tabular}
\end{center}
The tiny errors reflect the order \(4\) of the method (local error
\(-\frac1{240}h^6y^{(6)}\approx-4\times10^{-9}\) per step).
\end{solution}

%---------------------------------------------------------------------
\Q{Consistency and zero-stability for \(y''=f\)}{Short answer}{5}
\begin{qbox}
Verify that Störmer's method satisfies the consistency conditions for
\(y''=f\), and explain why the root condition for \(y''=f\) allows a
double root of \(\rho\) on the unit circle.
\end{qbox}

\begin{solution}
Störmer: \(\rho(\zeta)=(\zeta-1)^2\), \(\sigma(\zeta)=\zeta\).  Then
\(\rho(1)=0\), \(\rho'(1)=0\), \(\rho''(1)=2=2\sigma(1)\): the
consistency conditions \(C_0=C_1=C_2=0\) hold.

Consistency for \(y''=f\) \emph{forces} \(\zeta=1\) to be a double root
of \(\rho\) (since \(\rho(1)=\rho'(1)=0\)).  This mirrors the exact
problem: \(y''=0\) has the solutions \(1\) and \(x\), and the difference
equation \(\rho(E)y_n=0\) reproduces them as \(1\) and \(n=(x-a)/h\).
Because the starting values must also approximate the initial slope
(condition (ii) in the definition of convergence), the coefficient of
the component \(n\) is \(O(h)\), so this component stays bounded: it is
just the linear part \(x\,y'(a)\) of the solution.  A double root of
modulus one is therefore harmless, whereas multiplicity \(3\) would add
components like \(n^2\), which are not controlled in this way.
\end{solution}

%#####################################################################
\Needspace{32\baselineskip}
\section{Boundary Value Problems}
%#####################################################################

%---------------------------------------------------------------------
\Q{A Dirichlet problem by finite differences}{Numerical}{10}
\begin{qbox}
Solve \(y''-y=0\), \(y(0)=0\), \(y(1)=1\), with the central difference
scheme and \(h=\frac14\).  Solve the tridiagonal system by the Thomas
algorithm and compare with the exact solution
\(y=\sinh x/\sinh1\).
\end{qbox}

\begin{solution}
\step{The system}
\(\frac{y_{i-1}-2y_i+y_{i+1}}{h^2}-y_i=0\) gives
\(y_{i-1}-d\,y_i+y_{i+1}=0\) with \(d=2+h^2=2.0625\), for \(i=1,2,3\),
with \(y_0=0\), \(y_4=1\):
\[
  \begin{pmatrix}-d&1&0\\1&-d&1\\0&1&-d\end{pmatrix}
  \begin{pmatrix}y_1\\y_2\\y_3\end{pmatrix}
  =\begin{pmatrix}0\\0\\-1\end{pmatrix}.
\]
The matrix is strictly diagonally dominant (\(2.0625>2\)), so the
solution is unique.

\step{Forward sweep (express each unknown through the next)}
\begin{align*}
  \text{row 1:}&\quad y_1=\frac1d\,y_2=0.484848\,y_2,\\
  \text{row 2:}&\quad y_2=\frac1{d-0.484848}\,y_3=0.633854\,y_3,\\
  \text{row 3:}&\quad y_2-d\,y_3=-1\ \Rightarrow\
    y_3=\frac1{d-0.633854}=0.699963 .
\end{align*}

\step{Back substitution}
\(y_2=0.633854\times0.699963=0.443674\) and
\(y_1=0.484848\times0.443674=0.215115\).

\step{Comparison}
\begin{center}
\begin{tabular}{@{}cccc@{}}
\toprule
\(x\) & FD & exact & error\\
\midrule
0.25 & 0.215115 & 0.214952 & \(1.6\times10^{-4}\)\\
0.50 & 0.443674 & 0.443409 & \(2.6\times10^{-4}\)\\
0.75 & 0.699963 & 0.699724 & \(2.4\times10^{-4}\)\\
\bottomrule
\end{tabular}
\end{center}
The errors are \(O(h^2)\), as expected from the \(O(h^2)\) central
difference.
\end{solution}

%---------------------------------------------------------------------
\Q{A mixed boundary condition and tridiagonal form}{Derivation}{6}
\begin{qbox}
For \(y''=f(x)\) on \([0,1]\) with \(y'(0)-y(0)=1\), write the boundary
equation at \(x_0=0\) using the second-order one-sided formula for
\(y'(0)\).  Then eliminate \(y_2\) using the equation at \(x_1\) so that
the system becomes tridiagonal.
\end{qbox}

\begin{solution}
With \(y'(x_0)\approx\frac{-3y_0+4y_1-y_2}{2h}\), the boundary row is
\[
  -y_0+\frac{-3y_0+4y_1-y_2}{2h}=1 .
\]
It contains \(R_0y_2\) with \(R_0=-\frac1{2h}\).  The row at \(x_1\) is
\(\frac{y_0-2y_1+y_2}{h^2}=f_1\), with \(B_1=\frac1{h^2}\).  Form
row\(_0-\frac{R_0}{B_1}\)row\(_1\)\(=\)row\(_0+\frac h2\)row\(_1\):
\[
  -y_0+\frac{-3y_0+4y_1-y_2}{2h}+\frac{y_0-2y_1+y_2}{2h}=1+\frac h2f_1 .
\]
The \(y_2\) terms cancel, leaving
\[
  \Bigl(-1-\frac1h\Bigr)y_0+\frac1h\,y_1=1+\frac h2f_1,
\]
which has non-zero entries only in columns \(0\) and \(1\).
\end{solution}

%---------------------------------------------------------------------
\Q{Newton's method for a non-linear BVP}{Numerical}{6}
\begin{qbox}
For \(y''=y^2\), \(y(0)=y(1)=1\), with \(h=\frac12\), write the
discrete equation for the single unknown \(y_1\approx y(\frac12)\).
Perform one Newton step from \(y_1^{(0)}=1\) and compare with the exact
root of the discrete equation.
\end{qbox}

\begin{solution}
The scheme \(y_0-2y_1+y_2=h^2y_1^2\) with \(y_0=y_2=1\) gives
\[
  G(y_1)=2-2y_1-\tfrac14y_1^2=0,
  \qquad G'(y_1)=-2-\tfrac12y_1 .
\]
From \(y_1^{(0)}=1\): \(G(1)=-0.25\), \(G'(1)=-2.5\), so
\[
  y_1^{(1)}=1-\frac{G(1)}{G'(1)}=1-0.1=0.9 .
\]
The exact positive root is \(y_1=-4+2\sqrt6=0.898979\); one step already
has error \(1.0\times10^{-3}\), and the next step would square it
(quadratic convergence).  For larger \(N\) the Jacobian
\(J=A-h^2\,\mathrm{diag}(2y_i)\) is tridiagonal and each step is a
Thomas solve.
\end{solution}

%---------------------------------------------------------------------
\Q{Diagonal dominance and uniqueness}{Proof}{6}
\begin{qbox}
Prove that a strictly diagonally dominant matrix is non-singular.
Deduce that the central difference scheme for \(y''-q(x)y=f(x)\) with
Dirichlet conditions and \(q>0\) has a unique solution for every \(h\).
\end{qbox}

\begin{solution}
Suppose \(Ax=0\), \(x\ne0\), and pick \(i\) with
\(|x_i|=\max_j|x_j|>0\).  Row \(i\) gives
\(a_{ii}x_i=-\sum_{j\ne i}a_{ij}x_j\), so
\[
  |a_{ii}||x_i|\le\sum_{j\ne i}|a_{ij}||x_j|\le\Bigl(\sum_{j\ne i}|a_{ij}|\Bigr)|x_i|<|a_{ii}||x_i|,
\]
a contradiction.  So \(A\) is non-singular.

The scheme gives rows \(y_{i-1}-(2+h^2q_i)y_i+y_{i+1}=h^2f_i\).  The
diagonal entry has modulus \(2+h^2q_i>2\), and the off-diagonal entries
sum to at most \(2\) in modulus.  Hence the matrix is strictly
diagonally dominant and the solution is unique, for every \(h>0\).
\end{solution}

%---------------------------------------------------------------------
\Q{Reduction to homogeneous boundary conditions}{Short derivation}{5}
\begin{qbox}
For collocation applied to \(y''+p(x)y'+q(x)y=f(x)\) on \([0,1]\) with
\(y(0)=1\), \(y'(1)=2\), find the linear function \(d(x)\) used to make
the boundary conditions homogeneous, and the new right-hand side.  Why
does the reduction fail for the pure Neumann conditions
\(y'(0)=c_0\), \(y'(1)=c_1\)?
\end{qbox}

\begin{solution}
Here \(a_0=1,b_0=0\) at \(x=0\) and \(a_1=0,b_1=1\) at \(x=1\).  With
\(d=\alpha x+\beta\): \(d(0)=\beta=1\) and \(d'(1)=\alpha=2\), so
\(d(x)=2x+1\).  (The determinant condition
\((a_0b_1-a_1b_0)+a_0a_1(b-a)=1\ne0\) holds.)  Then \(\tilde y=y-d\)
satisfies \(\tilde y(0)=0\), \(\tilde y'(1)=0\) and
\[
  L\tilde y=f-L(d)=f(x)-2p(x)-(2x+1)q(x).
\]
For pure Neumann conditions \(a_0=a_1=0\), so the determinant
\((a_0b_1-a_1b_0)+a_0a_1(b-a)=0\): a linear \(d\) cannot match two
slope conditions unless \(c_0=c_1\), and then not uniquely.
\end{solution}

%#####################################################################
\Needspace{32\baselineskip}
\section{Classification and Difference Operators}
%#####################################################################

%---------------------------------------------------------------------
\Q{Classify and find characteristics}{Short answer}{6}
\begin{qbox}
Classify each equation and, where real, find the characteristic
directions \(dy/dx\).
\begin{enumerate}[label=(\alph*)]
  \item \(u_{xx}+4u_{xy}+4u_{yy}=0\);
  \item \(u_{xx}-2u_{xy}-3u_{yy}+u_x=0\);
  \item \(x\,u_{xx}+u_{yy}=0\).
\end{enumerate}
\end{qbox}

\begin{solution}
Compare with \(Au_{xx}+2Bu_{xy}+Cu_{yy}+\cdots\) and use
\(\frac{dy}{dx}=\frac{B\pm\sqrt{B^2-AC}}A\).
\step{(a)}
\(A=1\), \(B=2\), \(C=4\): \(B^2-AC=0\), \emph{parabolic}; one
(double) characteristic direction \(dy/dx=2\).
\step{(b)}
\(A=1\), \(B=-1\), \(C=-3\): \(B^2-AC=1+3=4>0\), \emph{hyperbolic};
\(dy/dx=-1\pm2\), i.e.\ \(1\) and \(-3\).  (The term \(u_x\) does not
affect the type.)
\step{(c)}
\(A=x\), \(B=0\), \(C=1\): \(B^2-AC=-x\).  Elliptic for \(x>0\),
parabolic on \(x=0\), hyperbolic for \(x<0\), where
\(dy/dx=\pm\sqrt{-x}/x=\mp1/\sqrt{-x}\).
\end{solution}

%---------------------------------------------------------------------
\Q{The central second difference}{Derivation}{6}
\begin{qbox}
Using \(E=e^{hD}\) and \(\delta=E^{1/2}-E^{-1/2}\), show that
\[
  \delta^2=h^2D^2+\frac{h^4}{12}D^4+\frac{h^6}{360}D^6+\cdots,
\]
and deduce the leading truncation error of
\(\frac{u_{m+1}-2u_m+u_{m-1}}{h^2}\approx u_{xx}\).  Then invert the
series to get the first two terms of \(D^2\) in powers of \(\delta\).
\end{qbox}

\begin{solution}
\(\delta=2\sinh\frac{hD}2\), so
\(\delta^2=4\sinh^2\frac{hD}2=2(\cosh hD-1)\).  With
\(\cosh t=1+\frac{t^2}2+\frac{t^4}{24}+\frac{t^6}{720}+\cdots\),
\[
  \delta^2=h^2D^2+\frac{h^4D^4}{12}+\frac{h^6D^6}{360}+\cdots .
\]
Hence
\(\frac{\delta^2u_m}{h^2}=u_{xx}+\frac{h^2}{12}u_{xxxx}+O(h^4)\): the
truncation error is \(\frac{h^2}{12}u_{xxxx}\).  Inverting,
\(h^2D^2=\delta^2-\frac{h^4D^4}{12}+\cdots=\delta^2-\frac{\delta^4}{12}+\cdots\)
(to this order \(h^4D^4=\delta^4\)), so
\(D^2=\frac1{h^2}\bigl(\delta^2-\frac1{12}\delta^4+\cdots\bigr)\), in
agreement with Note 4.
\end{solution}

%#####################################################################
\Needspace{32\baselineskip}
\section{Parabolic Equations: Schemes, Truncation Error, Stability}
%#####################################################################

Throughout, \(r=k/h^2\) and \(s=\sin\frac{\beta h}2\).

%---------------------------------------------------------------------
\Q{Two steps of the explicit scheme}{Numerical}{8}
\begin{qbox}
Solve \(u_t=u_{xx}\), \(u(x,0)=\sin\pi x\), \(u(0,t)=u(1,t)=0\), by the
explicit scheme with \(h=\frac14\), \(r=\frac14\).  Find \(U\) at
\(x=\frac12\) after one and two time steps and compare with the exact
solution \(e^{-\pi^2t}\sin\pi x\).
\end{qbox}

\begin{solution}
\(k=rh^2=\frac1{64}\).  Initially \(U^0=(0,\ \sin\frac\pi4,\ 1,\ \sin\frac{3\pi}4,\ 0)=(0,0.707107,1,0.707107,0)\).
\step{Step 1}
\(U_m^1=(1-2r)U_m^0+r(U_{m+1}^0+U_{m-1}^0)=\frac12U_m^0+\frac14(U_{m+1}^0+U_{m-1}^0)\):
\[
  U_2^1=\tfrac12(1)+\tfrac14(0.707107+0.707107)=0.853553,
  \quad
  U_1^1=\tfrac12(0.707107)+\tfrac14(1)=0.603553 .
\]
\step{Step 2}
\(U_2^2=\frac12(0.853553)+\frac14(2\times0.603553)=0.728553\).

\step{Why the numbers are so clean}
\(\sin\pi x_m\) is an eigenvector of the scheme: \(U_m^n=\xi^n\sin\pi x_m\)
with \(\xi=1-4r\sin^2\frac{\pi h}2=1-\sin^2\frac\pi8=\cos^2\frac\pi8=0.853553\).

\step{Comparison at \(x=\frac12\)}
\begin{center}
\begin{tabular}{@{}cccc@{}}
\toprule
\(t\) & explicit & exact & error\\
\midrule
\(1/64\) & 0.853553 & 0.857090 & \(3.5\times10^{-3}\)\\
\(2/64\) & 0.728553 & 0.734603 & \(6.0\times10^{-3}\)\\
\bottomrule
\end{tabular}
\end{center}
\end{solution}

%---------------------------------------------------------------------
\Q{Why ``unstable'' even when the solution decays}{Conceptual}{6}
\begin{qbox}
Repeat the previous problem with \(r=1\).  The computed solution still
decays.  Explain why the explicit scheme is nevertheless called unstable
for \(r=1\), and what would be seen in practice.
\end{qbox}

\begin{solution}
With \(h=\frac14\) the grid supports the modes \(\sin(s\pi x)\),
\(s=1,2,3\), with amplification factors
\(\xi_s=1-4r\sin^2\frac{s\pi}8\).  For \(r=1\):
\[
  \xi_1=1-4\sin^2\tfrac\pi8=0.414,\qquad
  \xi_2=1-4\sin^2\tfrac\pi4=-1,\qquad
  \xi_3=1-4\sin^2\tfrac{3\pi}8=-2.414 .
\]
The initial data contain only the mode \(s=1\), which decays.  But
rounding errors introduce small amounts of every mode, and the mode
\(s=3\) is multiplied by \(-2.414\) at every step: after \(20\) steps
it has grown by \((2.414)^{20}\approx4.5\times10^7\).  In practice a
saw-tooth oscillation appears and swamps the solution.  Stability is a
property of the scheme for \emph{all} data (all modes), not of one
solution; here \(\max_s|\xi_s|>1\), which violates \(r\le\frac12\).
\end{solution}

%---------------------------------------------------------------------
\Q{The \(\theta\)-method}{Stability analysis}{10}
\begin{qbox}
For \(0\le\theta\le1\), consider
\[
  U_m^{n+1}-U_m^n=r\bigl[\theta\,\delta^2U_m^{n+1}+(1-\theta)\,\delta^2U_m^n\bigr].
\]
Find the amplification factor and show that the scheme is stable for
all \(r>0\) if \(\theta\ge\frac12\), and for
\(r\le\frac1{2(1-2\theta)}\) if \(\theta<\frac12\).  Identify the
explicit, Crank--Nicolson and fully implicit schemes.
\end{qbox}

\begin{solution}
Substituting \(U_m^n=\xi^ne^{i\beta mh}\) and using
\(\delta^2e^{i\beta mh}=-4s^2e^{i\beta mh}\),
\[
  \xi-1=-4rs^2\bigl[\theta\xi+(1-\theta)\bigr]
  \ \Longrightarrow\
  \xi=\frac{1-4(1-\theta)rs^2}{1+4\theta rs^2}.
\]
Since \(rs^2\ge0\), \(\xi\le1\) always.  The condition \(\xi\ge-1\) is
\[
  1-4(1-\theta)rs^2\ge-1-4\theta rs^2
  \iff 4rs^2(1-2\theta)\le2 .
\]
If \(\theta\ge\frac12\) the left side is \(\le0\): stable for all \(r\).
If \(\theta<\frac12\), the worst case is \(s^2=1\), giving
\(r\le\frac1{2(1-2\theta)}\).

\(\theta=0\): explicit scheme, \(r\le\frac12\).
\(\theta=\frac12\): Crank--Nicolson, unconditionally stable.
\(\theta=1\): fully implicit, unconditionally stable.
(For example \(\theta=\frac14\) requires \(r\le1\).)
\end{solution}

%---------------------------------------------------------------------
\Q{A growing solution: \(u_t=u_{xx}+cu\)}{Stability analysis}{6}
\begin{qbox}
For \(u_t=u_{xx}+cu\) with \(c>0\), the explicit scheme is
\(U_m^{n+1}=U_m^n+r\delta^2U_m^n+ck\,U_m^n\).  Show that \(|\xi|\) can
exceed \(1\), but the scheme is still stable when \(r\le\frac12\).
\end{qbox}

\begin{solution}
\(\xi=1-4rs^2+ck\).  For \(s=0\), \(\xi=1+ck>1\).  If \(r\le\frac12\),
then \(1-4rs^2\in[-1,1]\), so
\[
  |\xi|\le1+ck\le e^{ck}
  \quad\Longrightarrow\quad
  |\xi|^n\le e^{cnk}=e^{cT}.
\]
The powers are bounded by a constant independent of \(h\) and \(k\)
(for \(nk\le T\)), which is stability in the Lax--Richtmyer sense
(\(K=1\), \(\beta=c\)).  The condition \(|\xi|\le1\) is only
appropriate when the exact solution does not grow; here the exact modes
grow like \(e^{ct}\), so \(|\xi|\le1+O(k)\) is the correct requirement.
\end{solution}

%---------------------------------------------------------------------
\Q{Truncation error of the fully implicit scheme}{Derivation}{6}
\begin{qbox}
Find the principal part of the local truncation error of
\(U_m^{n+1}-U_m^n=r\,\delta^2U_m^{n+1}\) for \(u_t=u_{xx}\), and state
its order of consistency.
\end{qbox}

\begin{solution}
\(T=u_m^{n+1}-u_m^n-r\delta^2u_m^{n+1}\).  Expand about
\((mh,(n+1)k)\):
\[
  u_m^n=u-ku_t+\tfrac12k^2u_{tt}-\cdots,
  \qquad
  r\delta^2u_m^{n+1}=k\,u_{xx}+\tfrac1{12}kh^2u_{xxxx}+\cdots .
\]
Hence
\[
  T=k(u_t-u_{xx})-\tfrac12k^2u_{tt}-\tfrac1{12}kh^2u_{xxxx}+\cdots
   =-\tfrac12k^2u_{tt}-\tfrac1{12}kh^2u_{xxxx}+\cdots,
\]
so \(T=O(k^2+kh^2)\) and \(T/k=O(k+h^2)\to0\): consistent, first order
in time and second order in space.  (Unlike the explicit scheme, no
choice of \(r\) cancels the principal part, because both terms have the
same sign when \(u_{tt}=u_{xxxx}\).)
\end{solution}

%---------------------------------------------------------------------
\Q{One Crank--Nicolson step}{Numerical}{8}
\begin{qbox}
For \(u_t=u_{xx}\), \(u(x,0)=\sin\pi x\), \(u(0,t)=u(1,t)=0\), take one
Crank--Nicolson step with \(h=\frac14\), \(r=1\).  Solve the system
and compare with the exact solution.
\end{qbox}

\begin{solution}
Crank--Nicolson,
\((1+r)U_m^{1}-\frac r2(U_{m+1}^1+U_{m-1}^1)=(1-r)U_m^0+\frac r2(U_{m+1}^0+U_{m-1}^0)\),
with \(r=1\) reads
\(2U_m^{1}-\frac12(U_{m+1}^1+U_{m-1}^1)=\frac12(U_{m+1}^0+U_{m-1}^0)\).
Right-hand sides: \(m=1\): \(\frac12(1+0)=0.5\); \(m=2\):
\(\frac12(0.707107\cdot2)=0.707107\); \(m=3\): \(0.5\).
\[
  \begin{pmatrix}2&-\frac12&0\\-\frac12&2&-\frac12\\0&-\frac12&2\end{pmatrix}
  \begin{pmatrix}U_1\\U_2\\U_3\end{pmatrix}
  =\begin{pmatrix}0.5\\0.707107\\0.5\end{pmatrix}.
\]
By symmetry \(U_1=U_3\).  Row 2: \(U_2=\frac12(0.707107+U_1)\).  Row 1:
\(2U_1-\frac14(0.707107+U_1)=0.5\Rightarrow\frac74U_1=0.676777\), so
\[
  U_1=U_3=0.386730,\qquad U_2=0.546918 .
\]
Check: \(U_2=\xi\) with
\(\xi=\frac{1-2r\sin^2(\pi/8)}{1+2r\sin^2(\pi/8)}=0.546918\).  Exact at
\(t=k=\frac1{16}\), \(x=\frac12\): \(e^{-\pi^2/16}=0.539641\); error
\(7.3\times10^{-3}\), even though \(r=1\) is twice the explicit limit.
\end{solution}

%---------------------------------------------------------------------
\Q{Matrix stability of the fully implicit scheme}{Proof}{6}
\begin{qbox}
Show by the matrix method that the fully implicit scheme with Dirichlet
conditions is unconditionally stable.
\end{qbox}

\begin{solution}
The scheme is \(A\mathbf U^{n+1}=\mathbf U^n+\mathbf b\) with
\(A=\mathrm{tridiag}(-r,\,1+2r,\,-r)\) of size \(N-1\).  By the
tridiagonal eigenvalue formula,
\[
  \lambda_s(A)=1+2r-2r\cos\frac{s\pi}N=1+4r\sin^2\frac{s\pi}{2N}\ \ge1 .
\]
\(A\) is real symmetric, so \(A^{-1}\) is symmetric with eigenvalues
\(1/\lambda_s(A)\in(0,1)\).  Hence
\(\|A^{-1}\|_2=\rho(A^{-1})<1\) for every \(r>0\), and
\(\|(A^{-1})^n\|_2\le1\): unconditionally stable.
\end{solution}

%---------------------------------------------------------------------
\Q{A Neumann condition by a fictitious point}{Derivation}{5}
\begin{qbox}
For the explicit scheme with \(u_x(0,t)=0\), derive the equation at
\(m=0\) using a fictitious point and central differences.  Compare its
accuracy with the condition \(U_0^n=U_1^n\).
\end{qbox}

\begin{solution}
Central difference: \(\frac{U_1^n-U_{-1}^n}{2h}=0\Rightarrow U_{-1}^n=U_1^n\).
The scheme at \(m=0\),
\(U_0^{n+1}=(1-2r)U_0^n+r(U_1^n+U_{-1}^n)\), becomes
\[
  U_0^{n+1}=(1-2r)U_0^n+2rU_1^n .
\]
The central approximation of \(u_x\) has error \(O(h^2)\), while
\(U_0=U_1\), i.e.\ \(\frac{U_1-U_0}h=0\), has error
\(\frac h2u_{xx}(0,t)=O(h)\).  So the fictitious-point treatment keeps
the second-order spatial accuracy of the interior scheme.  The
coefficients \(1-2r\) and \(2r\) are non-negative for \(r\le\frac12\),
so the maximum-principle stability argument still applies.
\end{solution}

%---------------------------------------------------------------------
\Q{Convergence of the explicit scheme}{Proof}{10}
\begin{qbox}
Prove that for \(u_t=u_{xx}\) on \([0,1]\) with Dirichlet conditions,
the explicit scheme converges for \(0<r\le\frac12\), with
\(\max_m|u_m^n-U_m^n|\le T\,A\,(k+h^2)\).
\end{qbox}

\begin{solution}
Let \(z_m^n=u_m^n-U_m^n\).  Subtracting the scheme from the equation
satisfied by the exact solution,
\[
  z_m^{n+1}=(1-2r)z_m^n+r(z_{m+1}^n+z_{m-1}^n)+T_m^n,
  \qquad z_0^n=z_N^n=0,
\]
with \(|T_m^n|\le A(k^2+kh^2)\), where
\(A=\max(\frac12\max|u_{tt}|,\frac1{12}\max|u_{xxxx}|)\).  For
\(r\le\frac12\) all coefficients are non-negative and sum to \(1\), so
with \(Z^{(n)}=\max_m|z_m^n|\),
\[
  Z^{(n+1)}\le Z^{(n)}+A(k^2+kh^2).
\]
Since \(Z^{(0)}=0\),
\(Z^{(n)}\le nA(k^2+kh^2)=nk\,A(k+h^2)=T\,A(k+h^2)\to0\).
\end{solution}

%---------------------------------------------------------------------
\Q{Time-step limits in one, two and three dimensions}{Short answer}{5}
\begin{qbox}
With \(h=0.1\), what is the largest stable time step for the explicit
scheme for the heat equation in 1D, 2D (five-point) and 3D?  What is it
for the 2D nine-point scheme, and for Crank--Nicolson or ADI?
\end{qbox}

\begin{solution}
In \(d\) dimensions the explicit scheme needs \(r\le\frac1{2d}\), i.e.
\(k\le\frac{h^2}{2d}\):
\[
  \text{1D: }k\le0.005,\qquad
  \text{2D: }k\le0.0025,\qquad
  \text{3D: }k\le0.00\overline{16}.
\]
The nine-point scheme has
\(\xi=(1-4rs_1^2)(1-4rs_2^2)\), stable for \(r\le\frac12\), so
\(k\le0.005\), twice the five-point limit.  Crank--Nicolson and the
Peaceman--Rachford ADI scheme are unconditionally stable: no limit (the
step is chosen for accuracy).
\end{solution}

%#####################################################################
\Needspace{32\baselineskip}
\section{Lecture 35: ADI, Three-Level and Du Fort--Frankel Schemes}
%#####################################################################

%---------------------------------------------------------------------
\Q{Peaceman--Rachford reproduces the factored scheme}{Proof}{6}
\begin{qbox}
Show that eliminating \(U^{n+1/2}\) from the Peaceman--Rachford
equations
\[
  \bigl(1-\tfrac r2\delta_{x_1}^2\bigr)U^{n+1/2}=\bigl(1+\tfrac r2\delta_{x_2}^2\bigr)U^n,
  \qquad
  \bigl(1-\tfrac r2\delta_{x_2}^2\bigr)U^{n+1}=\bigl(1+\tfrac r2\delta_{x_1}^2\bigr)U^{n+1/2}
\]
gives the two-dimensional Crank--Nicolson-type scheme.  What assumption
is needed?  Why is the split form preferred?
\end{qbox}

\begin{solution}
Apply \(\bigl(1-\frac r2\delta_{x_1}^2\bigr)\) to the second equation.
If \(\delta_{x_1}^2\) and \(\delta_{x_2}^2\) commute,
\[
  \bigl(1-\tfrac r2\delta_{x_1}^2\bigr)\bigl(1-\tfrac r2\delta_{x_2}^2\bigr)U^{n+1}
  =\bigl(1+\tfrac r2\delta_{x_1}^2\bigr)\bigl(1-\tfrac r2\delta_{x_1}^2\bigr)U^{n+1/2}
  =\bigl(1+\tfrac r2\delta_{x_1}^2\bigr)\bigl(1+\tfrac r2\delta_{x_2}^2\bigr)U^n .
\]
(The commutation holds for constant coefficients on a rectangle.)

The factored scheme couples all grid points at the new level through a
large banded system.  The split form needs only tridiagonal solves: in
the first half-step one tridiagonal system along each line in the
\(x_1\)-direction, in the second along each line in the
\(x_2\)-direction.  Each costs \(O(N)\) by the Thomas algorithm, so a
full step costs \(O(N^2)\) for \(N^2\) unknowns.
\end{solution}

%---------------------------------------------------------------------
\Q{Stability of the Peaceman--Rachford method}{Stability analysis}{8}
\begin{qbox}
Find the amplification factor of one full Peaceman--Rachford step and
show that the method is unconditionally stable.  Is each half-step
stable by itself?
\end{qbox}

\begin{solution}
Let \(U^n=\xi^ne^{i(\beta lh+\gamma mh)}\), \(s_1=\sin\frac{\beta h}2\),
\(s_2=\sin\frac{\gamma h}2\), and \(\xi^*\) the half-step factor.  Then
\[
  (1+2rs_1^2)\,\xi^*=1-2rs_2^2,
  \qquad
  (1+2rs_2^2)\,\xi=(1-2rs_1^2)\,\xi^*,
\]
so
\[
  \xi=\frac{1-2rs_1^2}{1+2rs_1^2}\cdot\frac{1-2rs_2^2}{1+2rs_2^2}.
\]
Each factor has the Crank--Nicolson form \(\frac{1-a}{1+a}\) with
\(a\ge0\), so \(|\xi|\le1\) for all \(r\), \(\beta\), \(\gamma\): the
method is unconditionally stable.

A half-step alone is not: \(\xi^*=\frac{1-2rs_2^2}{1+2rs_1^2}\) equals
\(1-2r\) when \(s_1=0\), \(s_2=1\), which has modulus greater than \(1\)
for \(r>1\).  Only the pair of half-steps is stable.
\end{solution}

%---------------------------------------------------------------------
\Q{Richardson's scheme is unstable: a determinant proof}{Proof}{6}
\begin{qbox}
The amplification matrix of the three-level scheme
\(U^{n+1}=2r\delta^2U^n+U^{n-1}\) is
\(G=\begin{pmatrix}-8rs^2&1\\1&0\end{pmatrix}\).  Without solving for the
eigenvalues, prove that \(G\) has an eigenvalue of modulus greater than
\(1\) whenever \(s\ne0\).
\end{qbox}

\begin{solution}
\(\det G=\lambda_1\lambda_2=-1\) and
\(\mathrm{trace}\,G=\lambda_1+\lambda_2=-8rs^2\).  The discriminant of
\(\lambda^2+8rs^2\lambda-1\) is \(64r^2s^4+4>0\), so the eigenvalues are
real and distinct.  Suppose both had modulus \(\le1\).  Since
\(|\lambda_1||\lambda_2|=1\), both have modulus exactly \(1\); being
real, they are \(\pm1\), and with product \(-1\) they are \(1\) and
\(-1\).  Then the trace is \(0\), contradicting
\(-8rs^2\ne0\).  Hence \(\max|\lambda_i|>1\) for every \(r>0\) and every
mode with \(s\ne0\): the scheme is unconditionally unstable.
\end{solution}

%---------------------------------------------------------------------
\Q{Stability of the Du Fort--Frankel scheme}{Stability analysis}{10}
\begin{qbox}
Prove that the Du Fort--Frankel scheme
\((1+2r)U_m^{n+1}=2r(U_{m+1}^n+U_{m-1}^n)+(1-2r)U_m^{n-1}\) is
unconditionally stable.
\end{qbox}

\begin{solution}
Substitute \(U_m^n=\xi^ne^{i\beta mh}\), \(\theta=\beta h\):
\[
  (1+2r)\xi^2-4r\cos\theta\,\xi-(1-2r)=0,
\]
whose roots are the eigenvalues of the \(2\times2\) amplification
matrix of the equivalent two-level system.  They are
\[
  \xi_{1,2}=\frac{2r\cos\theta\pm\sqrt{1-4r^2\sin^2\theta}}{1+2r}.
\]
\step{Case 1: \(4r^2\sin^2\theta\le1\) (real roots)}
\[
  |\xi_{1,2}|\le\frac{2r|\cos\theta|+\sqrt{1-4r^2\sin^2\theta}}{1+2r}
  \le\frac{2r+1}{1+2r}=1 .
\]
\step{Case 2: \(4r^2\sin^2\theta>1\) (complex conjugate roots)}
Then \(r>\frac12\), and
\(|\xi|^2=|\xi_1\xi_2|=\bigl|\frac{-(1-2r)}{1+2r}\bigr|=\frac{2r-1}{2r+1}<1\).

So \(|\xi_i|\le1\) for all \(r>0\) and all \(\theta\).  Roots of
modulus one occur only for \(\theta=0\) or \(\pi\), where they are
simple, so the powers of the amplification matrix stay bounded: the
scheme is unconditionally stable.
\end{solution}

%---------------------------------------------------------------------
\Q{Consistency of the Du Fort--Frankel scheme}{Derivation}{8}
\begin{qbox}
Show that the truncation error of the Du Fort--Frankel scheme (written
as an approximation to \(u_t=u_{xx}\)) is
\[
  T=\Bigl(\frac kh\Bigr)^2u_{tt}+\frac{k^2}6u_{ttt}-\frac{h^2}{12}u_{xxxx}+\cdots .
\]
When is the scheme consistent?  What happens if \(k/h=c\) is held fixed?
\end{qbox}

\begin{solution}
Write the scheme as
\[
  \frac{U_m^{n+1}-U_m^{n-1}}{2k}
  =\frac{U_{m+1}^n-U_m^{n+1}-U_m^{n-1}+U_{m-1}^n}{h^2}.
\]
Using
\(u^{n+1}-u^{n-1}=2ku_t+\frac{k^3}3u_{ttt}+\cdots\),
\(u^{n+1}+u^{n-1}=2u+k^2u_{tt}+\cdots\) and
\(u_{m+1}+u_{m-1}=2u+h^2u_{xx}+\frac{h^4}{12}u_{xxxx}+\cdots\),
\[
  T=u_t+\frac{k^2}6u_{ttt}-\Bigl[u_{xx}+\frac{h^2}{12}u_{xxxx}-\frac{k^2}{h^2}u_{tt}\Bigr]+\cdots
   =\Bigl(\frac kh\Bigr)^2u_{tt}+\frac{k^2}6u_{ttt}-\frac{h^2}{12}u_{xxxx}+\cdots,
\]
since \(u_t=u_{xx}\).  So \(T\to0\) if and only if \(k/h\to0\).  With
\(r=k/h^2\) fixed, \(k/h=rh\to0\), so the scheme is consistent.  But if
\(k/h=c\) is fixed, \(T\to c^2u_{tt}\ne0\): the scheme is then
consistent with the hyperbolic equation \(u_t+c^2u_{tt}=u_{xx}\), not
with the heat equation.  Unconditional stability does not mean that any
\(k\) gives a good answer.
\end{solution}

%---------------------------------------------------------------------
\Q{Du Fort--Frankel with a large time step}{Numerical}{6}
\begin{qbox}
For \(u_t=u_{xx}\), \(u(x,0)=\sin\pi x\), \(u(0,t)=u(1,t)=0\), take
\(h=\frac14\), \(r=1\), and the exact value at \(t=k\) as the second
starting level.  Compute \(U\) at \(x=\frac12\), \(t=2k\) by Du
Fort--Frankel and compare with the exact value.  Explain the size of
the error.
\end{qbox}

\begin{solution}
\(k=\frac1{16}\).  The data stay proportional to \(\sin\pi x\):
\(U_m^n=a_n\sin\pi x_m\), with \(U_{m+1}+U_{m-1}=2\cos\frac\pi4\,U_m\).
The scheme gives
\[
  3a_{n+1}=2\cdot2\cos\tfrac\pi4\,a_n-a_{n-1}=2.828427\,a_n-a_{n-1}.
\]
With \(a_0=1\), \(a_1=e^{-\pi^2/16}=0.539641\):
\[
  a_2=\frac{2.828427\times0.539641-1}{3}=0.175445,
\]
so \(U(\tfrac12,2k)=0.175445\), against the exact
\(e^{-2\pi^2/16}=0.291213\): an error of about \(40\%\).

The reason is the consistency error \((k/h)^2u_{tt}\).  Here
\(k/h=\frac14\), and for this mode \(u_{tt}=\pi^4u\approx97.4u\), so the
extra term is about \(6.1u\), comparable with \(u_t=-\pi^2u\approx-9.9u\).
The scheme is stable, but at this step size it is not accurate.
\end{solution}

%---------------------------------------------------------------------
\Q{Short answers on Lecture 35}{Short answer}{6}
\begin{qbox}
\begin{enumerate}[label=(\alph*)]
  \item Du Fort--Frankel contains \(U_m^{n+1}\) on both sides.  Is it
    explicit or implicit?
  \item What extra information does a three-level scheme need to start?
  \item Richardson's scheme has truncation error \(O(k^2+h^2)\), better
    than the explicit scheme.  Why is it useless?
  \item In ADI, what does ``alternating direction implicit'' refer to?
\end{enumerate}
\end{qbox}

\begin{solution}
\step{(a)}
Explicit: \(U_m^{n+1}\) is the only unknown at the new level in each
equation, so it is obtained directly by dividing by \(1+2r\); no system
has to be solved.

\step{(b)}
Two initial levels, \(U^0\) and \(U^1\).  \(U^1\) must be computed by a
two-level method (for example one explicit or Crank--Nicolson step, or
a Taylor expansion) of adequate accuracy.

\step{(c)}
It is unconditionally unstable (Q30): for every \(r>0\) some mode grows
like \(|\lambda_2|^n\) with \(|\lambda_2|>1\).  By the Lax equivalence
theorem it does not converge, whatever its truncation error.

\step{(d)}
Each half-step is implicit in one space direction and explicit in the
other, and the implicit direction alternates between \(x_1\) and
\(x_2\) from one half-step to the next.
\end{solution}

\end{document}
